\documentclass[conference]{IEEEtran}
\usepackage[T1]{fontenc}
\usepackage[utf8]{inputenc}
\usepackage{textcomp,amsmath,amssymb,graphicx,booktabs,tabularx,cite,url,xcolor}
\usepackage[hidelinks]{hyperref}
\newcolumntype{Y}{>{\raggedright\arraybackslash}X}
\newcommand{\todo}[1]{\textcolor{red}{[#1]}}   % every red note must be resolved before submission
\providecommand{\tightlist}{}
\begin{document}
\title{Know Your Agent: Verifiable Delegation for AI Agents That Pay}
\author{\IEEEauthorblockN{Swapnali Dashrath}
\IEEEauthorblockA{Independent AI Solution Architect and Researcher\\
Pune, India\\
ORCID 0009-0006-0366-1578}}
\maketitle
\begin{abstract}
AI agents can now pay on a person's behalf. In 2025 and 2026, card networks, payment providers and cloud platforms released products for agent payments. Most of them identify the agent and record what the user asked for, and some cap spending per session. We present Know Your Agent (KYA), a five-layer design in which a checkpoint outside the agent checks every payment against a signed mandate. To our knowledge, KYA adds three things that the public descriptions we reviewed do not combine: a speed breaker with a stated loss bound, a budget that spans all agents of one principal, and a tamper-evident trail that records refusals. We state the loss bound as a proposition with a proof. A prototype passes 23 functional tests and adds a median of about 0.6 ms per payment. Across seven failure scenarios, KYA blocked all seven, while a scoped bearer token blocked two and a token with a per-item limit blocked three. In a replay calibrated to a public runaway-agent incident, a per-item limit prevented no loss and the breaker held the loss to 200 US dollars. The breaker does not catch slow loops. A loop just under the speed limit lost 2,403 dollars in its first hour, and only the budget stopped it. The main open question is false positives on busy, legitimate agents, which we have not yet measured.
\end{abstract}
\begin{IEEEkeywords}
agentic AI, delegated authorization, agent identity, runtime verification, policy enforcement, digital payments
\end{IEEEkeywords}

\section{Introduction}

Agents with payment authority are moving into production. Visa's Trusted Agent Protocol lets merchants recognize registered agents \cite{b1}. Mastercard Agent Pay issues payment tokens bound to registered agents \cite{b2}. Google's Agent Payments Protocol (AP2) records signed mandates of what the user wanted \cite{b3}. Amazon made AgentCore Payments generally available on August 18, 2026. Its session spending limits are enforced at the infrastructure layer, and a payment that would exceed the limit is denied \cite{b4}. In India, NPCI is reported to be building a Unified Agent Protocol for UPI. At Global Fintech Fest in September 2026, its chairman said that NPCI is examining protocols to identify and authorise agents on UPI \cite{b5}. We found no public specification.

The largest loss does not always come from an attacker. In September 2026, Mandiant reported an accounting agent that entered a recursive loop after a corrupted value broke one of its tools. The agent made more than 15,000 high-cost API calls in under an hour, which Mandiant put at about 50,000 US dollars in cloud charges \cite{b6,b7}. These figures come from the vendor and have not been independently verified \cite{b8}. Nobody attacked the agent, and nothing stopped it. The charges were metered cloud usage, not card payments, but the shape is the same as a payment loop: one small, individually valid charge repeated thousands of times. An agent that holds a payment token can fail in the same way.

We argue that this is a software engineering problem before it is a payments problem. Where should a check run so that the component it guards cannot skip it? What contract binds an autonomous component to the system it acts on? What evidence does the system leave behind? Current products answer part of this. They identify the agent, sign the request, and in some cases cap spending. None of the public descriptions we reviewed checks each payment against a signed mandate together with a running budget and a speed limit, and none keeps a record of refusals that a later dispute could use.

This paper makes three contributions:

\begin{itemize}
\item
  Know Your Agent (KYA), a five-layer design with a threat model of six failure patterns, five stated invariants, and a loss bound for runaway loops (Section III).
\item
  An open-source prototype with functional tests, latency measurements, a comparison with two common controls, and a replay of the public incident that includes the slow loops the breaker misses (Section IV) \cite{b31}.
\item
  A research agenda of four open questions for the software engineering community (Section V).
\end{itemize}

\section{Background and Related Work}

\subsection{Why agents break delegation}

Three properties of LLM agents change the problem. First, the agent chooses its actions at run time, so the set of possible actions is not known when a token is issued. Second, the agent reads untrusted content, and hidden instructions in that content can change what it does \cite{b9}. OWASP lists this risk together with excessive agency among the top risks for LLM applications \cite{b10}. A hijacked agent still holds a valid token, so identity checks still pass. Third, agents hand work to other agents while the task runs, so the chain of who acts for whom is built at run time. The combined result is a stream of actions that are each allowed on their own, such as one more small payment repeated thousands of times.

\subsection{Related work}

\emph{Delegation tokens.} OAuth 2.0 fixes a token's scope when it is issued \cite{b11}, and token exchange lets one party act for another while recording the acting party \cite{b12}. Rich Authorization Requests add structured, type-specific permissions to the token request \cite{b13}. GNAP generalizes these ideas with key-bound tokens and explicit grant revocation \cite{b14}. All of them answer whether a client may call an API. None tracks how much has been spent across calls or how quickly the calls arrive.

\emph{Attenuation.} Macaroons let a token holder add restrictions offline, without asking the issuer \cite{b15}. Biscuit does the same with public-key signatures and a logic language \cite{b16}, and UCAN defines capability chains for decentralized delegation \cite{b17}. Our sub-agent rule belongs to this family. These formats restrict what a token may do, but they do not count what has already been spent.

\emph{Workload identity.} SPIFFE gives each workload a verifiable identity \cite{b18}. It proves which software is calling, not whether that software was steered by injected text.

\emph{Policy engines.} OPA \cite{b19} and Cedar \cite{b20} evaluate policies at decision time, and a checkpoint can use either for static rules. Neither keeps spend counters or sliding-window counts unless the surrounding system adds them.

\emph{Runtime enforcement.} CaMeL extracts the control flow from the user's request, keeps untrusted data from changing that flow, and enforces policies in an interpreter without changing the underlying model \cite{b21}. KYA shares CaMeL's idea of enforcement outside the model, but it protects payment state rather than data flow. Reference monitors \cite{b22} and complete mediation \cite{b23} give the principle that every action passes through the check, and the circuit breaker pattern \cite{b24} gives the stop rule.

\emph{Agent payment products.} Visa, Mastercard and Google sign or bind agent requests and mandates \cite{b1,b2,b3}. AgentCore Payments enforces session-level limits and states that limits must be enforced outside the model \cite{b4,b25}. KYA builds on both ideas. It adds three things that we could not find in the public documentation we reviewed: a speed breaker with a stated loss bound, a budget aggregated across all agents of one principal, and a hash-chained trail that includes refusals.

\emph{Agent security and governance.} Indirect prompt injection \cite{b9} and agent benchmarks \cite{b26} define the threat and evaluation context for our work. Research on agent credentials \cite{b27} and on visibility into agent activity \cite{b28} is the closest prior work. Zero trust architecture \cite{b29} supports the principle that no component should be trusted only because it sits inside the system.

\section{The Know Your Agent Design}

\subsection{Threat model}

A principal lets an agent pay for a task. The agent runs an LLM planner with tools and memory, and it reads untrusted content. We trust the checkpoint, the registry and the principal's key store. We do not trust the agent, which may loop, be hijacked or be impersonated. Table I lists six failure patterns drawn from public incidents, protocol documents and security research.

\begin{table*}[t]
\caption{Six failure patterns and the layer that addresses each}
\centering\footnotesize
\begin{tabularx}{\textwidth}{@{}l>{\raggedright\arraybackslash}p{0.17\textwidth}>{\raggedright\arraybackslash}p{0.12\textwidth}Y>{\raggedright\arraybackslash}p{0.27\textwidth}@{}}
\toprule
\textbf{ID} & \textbf{Pattern} & \textbf{Attacker} & \textbf{What happens} & \textbf{Stopped by} \\
\midrule
T1 & Agent that will not stop & None & A loop repeats a paid action thousands of times \cite{b6} & L3 speed breaker and budget \\
T2 & Poisoned page & Content author & Hidden text changes the payee or amount \cite{b9} & L3 mandate scope (partial, see below) \\
T3 & Fake agent & Token thief & A bot presents a stolen token or claims a trusted agent's identity & L1 bound token and agent signature \\
T4 & Stale permission & None & The agent keeps paying after the mandate ends or was edited & L2 signature and expiry \\
T5 & Agents that team up & Colluding agents & Several agents split work to stay under each one's limit & L3 per-principal budget (partial) \\
T6 & Dispute nobody can settle & None & After a complaint, no one can prove who approved what & L4 decision trail \\
\bottomrule
\end{tabularx}
\end{table*}

KYA does not stop every injection. An injected instruction that keeps the payee, amount, budget and speed inside the mandate gets through, because from the checkpoint's view that payment looks like a legitimate purchase. KYA limits how much such a payment can take and how often it can happen, and the trail records it afterward. It does not detect the intent behind the payment. That would need a separate check on the agent's purpose, which we leave to future work.

\subsection{The five layers}

Fig. 1 shows one payment passing through KYA. Each layer has one job and passes evidence to the next.

\begin{figure}[t]
\centering
\includegraphics[width=\columnwidth]{kya_fig1.pdf}
\caption{One agent payment through KYA. Every payment passes through the L3 checkpoint. The dashed line is a step-up request, sent only when a rule fires.}
\label{fig1}
\end{figure}

\emph{L1: Agent identity.} Each agent has its own Ed25519 key pair. The registry issues a payment token that names one agent and carries the registry's signature. The agent signs each proposal over its agent ID, the payee, the amount and the token ID. The checkpoint verifies three things: the registry signed the token, the token names this agent, and the agent signed this proposal. A stolen token is useless without the agent's key, and an agent never borrows the user's identity.

\emph{L2: Signed mandate.} The principal's permission is written as a mandate and signed with the principal's key. A mandate names the principal, the agent and the allowed payees. It sets a per-item limit, a total budget, an expiry, and a speed limit of at most $n$ payments in a window of length $W$. It may also set a step-up threshold. The mandate ID is a hash of the signed content, so any later edit produces a mandate whose signature fails. The checkpoint refuses any mandate that leaves the per-item limit, budget, speed limit or expiry unset.

\emph{L3: Policy checkpoint.} The checkpoint sits between the agent and the payment network, and every payment must pass through it. For each proposal it checks, in order: identity (L1), mandate signature and completeness, suspension state, expiry and payee, per-item limit, mandate budget, the principal's budget across all agents, and the speed window. The speed window counts executed payments under one mandate. If a proposal passes every earlier check and the window already holds $n$ payments, the checkpoint refuses it and suspends the mandate. Proposals that fail an earlier check do not count toward the window, so forged or out-of-scope proposals cannot suspend a legitimate mandate. Repeated refusals are instead raised as an alert from the trail. Only a proposal that passes all of these checks reaches the bank's existing fraud engine, which receives agent context: agent ID, mandate ID, budget left, recent activity and past denials. KYA does not replace fraud scoring, and a low fraud score never overrides a scope failure.

\emph{L4: Decision trail.} Every decision is appended to a trail: allow, deny, flag, and any human answer. Each record holds the mandate ID, principal, agent, payee, amount, result, fraud score, the hash of the previous record, and a signature from the trail's key. Editing any record breaks the chain from that point on. Refusals are recorded as well, because a refusal is often the evidence a dispute needs. This follows the call for agent activity logs in \cite{b28}. A hash chain alone does not protect against an operator who rebuilds the whole chain, so a deployment should publish the latest record hash to an outside party at regular intervals.

\emph{L5: Human step-up.} Asking the human every time would make the agent useless. KYA asks only when a rule fires: a tripped breaker, a high fraud score, or an amount above the mandate's step-up threshold. It also notifies the human, without blocking, when it sees a run of payments just under the per-item limit. The answer is written to the trail.

\subsection{Concurrency, refunds and failures}

A production checkpoint runs on several replicas, so counters must be updated atomically. The design uses reserve-then-commit. Before the checkpoint forwards a payment, it reserves the amount against the mandate budget, the principal budget and a slot in the speed window, in one conditional update that succeeds only if the new totals stay within limits. A successful payment turns the reservation into a commit, and a failed or timed-out payment releases it. Each reservation carries the proposal's idempotency key, so a retried proposal cannot reserve twice. A reservation that outlives its timeout is released only after the network confirms that no payment exists. A refund reduces committed spend only when the network confirms it against a known payment ID, and committed spend never falls below zero. Without atomic reservation, each of $k$ unsynchronized replicas could admit $n$ payments before seeing the others, which weakens the bound in Section IV to $k \cdot n \cdot c$.

\subsection{Invariants}

We state five invariants that the checkpoint must keep under any sequence of proposals, failures and replicas.

\begin{itemize}
\item
  \emph{I1 (budget).} Committed plus reserved spend under a mandate never exceeds its total budget.
\item
  \emph{I2 (principal).} Committed plus reserved spend across all mandates of one principal never exceeds the principal's cap.
\item
  \emph{I3 (scope).} No payment executes unless its mandate signature verifies, the mandate is unexpired and not suspended, and the payee is listed.
\item
  \emph{I4 (breaker).} After a suspension, no payment executes until a human reinstates the mandate.
\item
  \emph{I5 (trail).} Any edit to a stored record breaks verification of the chain from that record onward.
\end{itemize}

We plan to check these with property-based state-machine tests that generate random sequences of proposals, failures, refunds and replica interleavings \cite{b30}.

\subsection{What the prototype leaves out}

The prototype runs in one process and implements the five layers as described in Section III-B. Six parts of the design are not in it, and we do not claim them in the evaluation.

\begin{itemize}
\item
  \emph{Replay protection.} The design adds the mandate ID, a nonce and a timestamp to the signed proposal, and the checkpoint rejects stale timestamps and repeated nonces. In the prototype a captured proposal can be resent.
\item
  \emph{Reserve-then-commit} across replicas (Section III-C).
\item
  \emph{Reinstatement and revocation.} The prototype suspends a mandate automatically. It has no call for a human to reinstate a suspended mandate or to revoke an agent or mandate at any time.
\item
  \emph{The sub-agent rule.} A sub-agent gets its own identity and mandate. Its mandate may not allow more payees, a higher limit, a larger budget or a faster pace than its parent's, and its spend counts against the same principal.
\item
  \emph{External anchoring} of the trail (L4).
\item
  \emph{The invariant tests} (Section III-D).
\end{itemize}

\section{Evaluation}

We asked four questions. Does each layer do what we claim? What does the checkpoint cost per payment? Which failures does KYA block that common controls do not? Does the design change the outcome of a real incident?

\subsection{Setup}

The reference implementation is about 850 lines of Python. It uses the cryptography library for Ed25519 signatures. It includes the five layers and an adapter to a fraud engine. A transparent reference engine stands in for a bank's real one, and a simulated human makes approval decisions. All experiments run this code in one process with synthetic workloads. The latency run used an Intel Core i5-4300M CPU (2.60 GHz, 2 cores) with Windows 10 (build 19045) and Python 3.14.3.

\subsection{Functional tests}

We wrote 23 tests that cover each layer and each failure pattern, and all 23 passed. The groups are identity (2 tests), mandate (3), checkpoint (5), human step-up (3), decision trail (2), and fraud-engine integration (8). Three examples show the behavior. A stolen token used with the thief's key is denied with a bad-agent-signature error. A loop of 80 calls in 18 seconds allows 60 calls and then suspends the mandate. Editing record 12 of the trail is detected at record 12. Two further results show that the layers depend on each other. A low fraud score did not let a payment to an unlisted payee through, because scope is checked first. In the other, two agents shared one principal with a 150-dollar cap and took turns making 50-dollar payments. The first three were allowed and the fourth was refused, although each agent had used at most 100 dollars of its own 5,000-dollar budget.

\subsection{Latency}

Over 3,000 checks on the allow path, a full check took a median of 632 \textmu s (99th percentile 888 \textmu s). With the fraud engine, monitoring and case handling attached, the median was 654 \textmu s (99th percentile 796 \textmu s). The gap between the two runs is within run-to-run variation. Each check verifies three Ed25519 signatures and creates one. The prototype is unoptimized and single-threaded, and the numbers exclude network time to a real bank. They show that the check itself adds little delay.

\subsection{Comparison with common controls}

Table II compares three setups on seven scenarios. The first baseline is a signed bearer token with a scope and an expiry, as in OAuth. The second adds a per-item limit, as in most payment products. The third is the KYA prototype. We implemented both baselines as a small token gateway and ran the same scripted scenarios against all three.

\begin{table}[t]
\caption{Outcome by scenario and setup}
\centering\footnotesize
\begin{tabularx}{\columnwidth}{@{}>{\raggedright\arraybackslash}p{0.30\columnwidth}YYY@{}}
\toprule
\textbf{Scenario} & \textbf{Token only} & \textbf{Token + per-item limit} & \textbf{KYA} \\
\midrule
Injected payee & Allowed & Allowed & Blocked \\
Stolen token & Allowed & Allowed & Blocked \\
Expired mandate & Blocked & Blocked & Blocked \\
Edited mandate & Blocked & Blocked & Blocked \\
Runaway loop, 80 calls & 80 allowed & 80 allowed & 60 allowed, then suspended \\
Two agents, 150-dollar cap & 200 dollars paid & 200 dollars paid & 150 dollars paid, then refused \\
Single oversized payment & Allowed & Blocked & Blocked \\
\bottomrule
\end{tabularx}
\end{table}

Two rows show no difference, and we report them for that reason. A signed token already carries an expiry and already fails verification when it is edited. The differences are in the other five rows. Two points limit what the table shows. The baseline token's scope is the payment API, not a payee list. A token with payee-level permissions, as Rich Authorization Requests allow \cite{b13}, would block the injected payee. The baselines are also our own minimal implementations, not the commercial products, which we could not run. The sub-agent scenario is not evaluated because the rule is not implemented.

\subsection{Replay and the loss bound}

Consider a mandate that allows at most $n$ payments in any window of length $W$, with each payment costing at most $c$, and a total budget $B$.

\textit{Proposition 1.} Assume I1 and I4 hold. If $n{+}1$ consecutive proposals that pass every other check arrive within a span shorter than $W$, then at most $n$ of them execute and the mandate is suspended at the last one. The loss from that burst before a human acts is at most $n \cdot c$. Independently, the total loss under the mandate is at most $B$.

\textit{Proof.} Let $t_1 < \cdots < t_{n+1}$ be the times of the proposals, with $t_{n+1} - t_1 < W$. Suppose the first $n$ all execute. When proposal $n{+}1$ arrives, each of them is younger than $W$, so the window holds at least $n$ executed payments. The checkpoint therefore refuses proposal $n{+}1$ and suspends the mandate. If the mandate was suspended earlier, fewer than $n$ executed. In both cases at most $n$ payments execute, each costing at most $c$. By I4, no later proposal executes until a human acts. The bound $B$ follows from I1. \hfill$\blacksquare$

For a loop that proposes at a constant rate $r$, the $n{+}1$ proposals span $n/r$ seconds, so the condition holds exactly when $r > n/W$. Loops above that rate are bounded by $n \cdot c$. Slower loops never trip the breaker, and for them only the budget $B$ bounds the loss, which grows as $r \cdot c \cdot t$ until $B$ is reached. The bound is per mandate. If $k$ agents of one principal loop at once, each under its own mandate, the loss is at most $k \cdot n \cdot c$, or the principal's cap if that is lower. The same factor applies to $k$ replicas without atomic reservation (Section III-C).

\begin{table}[t]
\caption{One-hour runaway replay. synthetic workload calibrated to a public incident \cite{b6,b7}}
\centering\footnotesize
\begin{tabularx}{\columnwidth}{@{}>{\raggedright\arraybackslash}p{0.52\columnwidth}YY@{}}
\toprule
\textbf{Setting} & \textbf{Loss (USD)} & \textbf{Loop stopped at} \\
\midrule
No controls & 54,543 & Never \\
100-dollar per-item limit only & 54,543 & Never \\
5,000-dollar total budget only & 5,000 & 330 s \\
Full KYA (budget plus 60 payments per 300 s) & 200 & About 13 s \\
\bottomrule
\end{tabularx}
\end{table}

The replay instantiates the proposition. The agent proposed about 4.5 times per second at about 3.33 dollars per call, so $n \cdot c = 60 \times 3.33 \approx 200$ dollars. The breaker fired at the 61st proposal, about 13 seconds into the run. The per-item limit prevented nothing, since each call cost far less than 100 dollars. A 5,000-dollar budget on its own capped the loss at that amount, but only after 330 seconds. We also varied the loop rate from 0.5 to 20 calls per second, and the loss stayed at about 200 dollars in every run.

Table IV shows what happens around and below the threshold $n/W = 0.2$ calls per second, under the full KYA mandate.

\begin{table}[t]
\caption{Loops near and below the breaker threshold (full kya mandate)}
\centering\footnotesize
\begin{tabularx}{\columnwidth}{@{}YYYY@{}}
\toprule
\textbf{Loop rate (calls/s)} & \textbf{Loss in first hour (USD)} & \textbf{Stopped by} & \textbf{Stopped at} \\
\midrule
0.05 & 603 & Nothing within 6 hours & Not stopped \\
0.10 & 1,203 & Budget & 15,000 s \\
0.19 & 2,283 & Budget & 7,895 s \\
0.20 & 2,403 & Budget & 7,500 s \\
0.21 & 200 & Breaker & 286 s \\
0.25 & 200 & Breaker & 240 s \\
\bottomrule
\end{tabularx}
\end{table}

The breaker fires only above the threshold, as the proposition predicts. A loop at exactly 0.2 calls per second ran for more than two hours and lost the full 5,000-dollar budget. Its first-hour loss was twelve times that of the fast loop. Under this mandate a slow loop is therefore more costly than a fast one, and the budget is the only control that stops it. These runs check that the code matches the bound. They are not evidence that KYA would stop a real incident, because real loops differ in cost distribution and in how steadily they propose.

\section{Limits and Research Agenda}

KYA is tested as a single-process prototype with a simulated human, a stand-in fraud engine and scripted workloads. It is not a live deployment, and Section III-E lists the parts of the design that the prototype leaves out. We compared commercial products from their public descriptions, not by running them. The design also assumes the checkpoint cannot be bypassed. Making that true is a deployment problem, because routing and token design must force every payment through the checkpoint.

Three limits follow from the results. First, the speed limit is per mandate. Only the budget is shared across a principal's agents, so collusion (T5) is handled by the principal's cap alone, and agents that coordinate through shared channels can spread work so that no single limit sees it. Second, the breaker misses slow loops (Table IV). A second, longer window or a limit derived from the task would narrow that gap. Third, a speed limit can pause legitimate work, such as a procurement agent that places many small orders on a busy day. A breaker that pauses a good task costs minutes, while a missing breaker can cost the whole credit line, but we have not measured how often such pauses happen.

Two studies are planned and have not been run. The first is a benign-workload study with a procurement agent, a refund-processing agent and a per-call research agent, each running a day of synthetic traffic with peaks. For each speed limit we will report how often the breaker suspends a mandate that should not be suspended, how many payments are paused, and for how long. We will compare fixed limits with limits derived from a warm-up profile. The second is a study with a real LLM agent in a sandbox, running shopping tasks in which some pages carry injected instructions in the style of AgentDojo \cite{b26}. We will report which attacks reach the payment step, and the benign task completion rate, with and without KYA.

For the software engineering community, we see four open questions. How do we turn a user's natural-language request into a complete, checkable mandate, and detect when a request is too vague to bound? How do we test an agent and its checkpoint against injected content at scale? How do we detect several agents that each stay within their limits but together do not? How do we derive speed and budget limits from task profiles, so that the breaker stops a loop but not a busy, legitimate agent?

\section{Conclusion}

AI agents can now hold a payment token and press Buy. Today's products can name an agent, sign the request, and in some cases cap a session's spend. KYA checks every payment against a signed mandate and adds a speed breaker with a stated loss bound, a budget that spans all agents of a principal, and a trail that keeps refusals. In our prototype, KYA blocked seven of seven failure scenarios, against two and three for the token baselines. A per-item limit, the most familiar control, prevented none of the loss in a replayed incident, while the breaker held it to the bound in Proposition 1. The same experiments show where the design stops: a loop slower than the speed limit is caught only by the budget. The next steps are a false-positive study on benign workloads, a live LLM agent study, and property-based tests of the five invariants.

\section*{Acknowledgment}

\textbf{Use of AI tools.} The author used Claude, an AI assistant from Anthropic, to help draft and edit text in all sections of this paper and to write the comparison script used in Section IV-D. The author designed the architecture, wrote the reference implementation, and is responsible for all content, including the experimental results.

\textbf{Conflict of interest.} The author previously worked at Visa Inc.~and completed Mastercard-related client engagements through prior employers. This paper uses only public information about both companies. The author has no current financial or advisory relationship with either.

\textbf{Code availability.} The reference implementation, tests and comparison script are open source \cite{b31}. \todo{ADD kya\_table3.py TO THE REPO, UPDATE THE README, TAG A NEW RELEASE AND UPDATE THE VERSION IN REF. 31}

\begin{thebibliography}{00}
\bibitem{b1} Visa Inc., ``Visa introduces Trusted Agent Protocol: An ecosystem-led framework for AI commerce,'' press release, Oct.~14, 2025. {[}Online{]}. Available: https://usa.visa.com/about-visa/newsroom/press-releases.releaseId.21716.html (accessed Nov.~12, 2026).

\bibitem{b2} Mastercard, ``Mastercard unveils Agent Pay, pioneering agentic payments technology to power commerce in the age of AI,'' press release, Apr.~29, 2025. {[}Online{]}. Available: https://www.mastercard.com/us/en/news-and-trends/press/2025/april/mastercard-unveils-agent-pay-pioneering-agentic-payments-technology-to-power-commerce-in-the-age-of-ai.html (accessed Nov.~12, 2026).

\bibitem{b3} S. Parikh and R. Surapaneni, ``Powering AI commerce with the new Agent Payments Protocol (AP2),'' Google Cloud Blog, Sep.~16, 2025. {[}Online{]}. Available: https://cloud.google.com/blog/products/ai-machine-learning/announcing-agents-to-payments-ap2-protocol (accessed Nov.~12, 2026).

\bibitem{b4} Amazon Web Services, ``AgentCore payments is now generally available in Amazon Bedrock AgentCore,'' What's New, Aug.~18, 2026. {[}Online{]}. Available: https://aws.amazon.com/about-aws/whats-new/2026/08/bedrock-agentcore-payments-ga/ (accessed Nov.~12, 2026).

\bibitem{b5} MediaNama, ``AI agents can read intent, not approve payments: NPCI chairman at GFF 2026,'' Sep.~2026. [Online]. Available: https://www.medianama.com/2026/09/223-npci-ai-agents-upi-payments/ (accessed Nov.~13, 2026).

\bibitem{b6} Mandiant, ``AI risk and resilience in 2026: A Mandiant special report,'' Google Cloud, Sep.~2026. [Online]. Available: https://cloud.google.com/security/resources/ai-risk-and-resilience-2026 (accessed Nov.~13, 2026).

\bibitem{b7} J. Pogorelec, ``One runaway AI agent racked up a \$50,000 cloud bill,'' Help Net Security, Sep.~16, 2026. {[}Online{]}. Available: https://www.helpnetsecurity.com/2026/09/16/google-mandiant-enterprise-ai-security-risks-report/ (accessed Nov.~12, 2026).

\bibitem{b8} Cloud Security Alliance, ``Frontier Ready Daily: 16 September 2026,'' Sep.~16, 2026. {[}Online{]}. Available: https://labs.cloudsecurityalliance.org/research/frontier-ready-daily-20260916/ (accessed Nov.~12, 2026).

\bibitem{b9} K. Greshake et al., ``Not what you've signed up for: Compromising real-world LLM-integrated applications with indirect prompt injection,'' in Proc. 16th ACM Workshop on Artificial Intelligence and Security (AISec), 2023, pp.~79--90.

\bibitem{b10} OWASP Foundation, ``OWASP Top 10 for LLM Applications 2025,'' Nov.~17, 2024. {[}Online{]}. Available: https://genai.owasp.org/resource/owasp-top-10-for-llm-applications-2025/ (accessed Nov.~12, 2026).

\bibitem{b11} D. Hardt, Ed., ``The OAuth 2.0 Authorization Framework,'' IETF RFC 6749, Oct.~2012.

\bibitem{b12} M. Jones et al., ``OAuth 2.0 Token Exchange,'' IETF RFC 8693, Jan.~2020.

\bibitem{b13} T. Lodderstedt, J. Richer, and B. Campbell, ``OAuth 2.0 Rich Authorization Requests,'' IETF RFC 9396, May 2023.

\bibitem{b14} J. Richer, Ed., and F. Imbault, ``Grant Negotiation and Authorization Protocol (GNAP),'' IETF RFC 9635, Oct.~2024.

\bibitem{b15} A. Birgisson et al., ``Macaroons: Cookies with contextual caveats for decentralized authorization in the cloud,'' in Proc. Network and Distributed System Security Symposium (NDSS), 2014.

\bibitem{b16} G. Couprie, ``Biscuit, the foundation for your authorization systems,'' Clever Cloud Engineering Blog, Apr.~12, 2021. {[}Online{]}. Available: https://www.clever.cloud/blog/engineering/2021/04/12/introduction-to-biscuit/ (accessed Nov.~12, 2026).

\bibitem{b17} UCAN Working Group, ``UCAN Specification.'' {[}Online{]}. Available: https://github.com/ucan-wg/spec (accessed Nov.~12, 2026).

\bibitem{b18} SPIFFE Project, ``SPIFFE ID and Verifiable Identity Document,'' specification. {[}Online{]}. Available: https://spiffe.io/docs/latest/spiffe-specs/spiffe-id/ (accessed Nov.~12, 2026).

\bibitem{b19} Open Policy Agent. {[}Online{]}. Available: https://www.openpolicyagent.org/ (accessed Nov.~12, 2026).

\bibitem{b20} J. W. Cutler et al., ``Cedar: A new language for expressive, fast, safe, and analyzable authorization,'' Proc. ACM Program. Lang., vol.~8, no. OOPSLA1, pp.~670--697, Apr.~2024, doi: 10.1145/3649835.

\bibitem{b21} E. Debenedetti et al., ``Defeating prompt injections by design,'' arXiv:2503.18813, 2025. {[}Online{]}. Available: https://arxiv.org/abs/2503.18813

\bibitem{b22} J. P. Anderson, ``Computer security technology planning study,'' Air Force Electronic Systems Division, Hanscom AFB, MA, USA, Tech. Rep.~ESD-TR-73-51, vols. I and II, Oct.~1972.

\bibitem{b23} J. H. Saltzer and M. D. Schroeder, ``The protection of information in computer systems,'' Proc. IEEE, vol.~63, no. 9, pp.~1278--1308, Sep.~1975.

\bibitem{b24} M. T. Nygard, \emph{Release It! Design and Deploy Production-Ready Software}, 2nd ed.~Raleigh, NC, USA: Pragmatic Bookshelf, 2018.

\bibitem{b25} Amazon Web Services, ``Enable safe agentic payments with built-in guardrails using Amazon Bedrock AgentCore payments,'' AWS Machine Learning Blog. {[}Online{]}. Available: https://aws.amazon.com/blogs/machine-learning/enable-safe-agentic-payments-with-built-in-guardrails-using-amazon-bedrock-agentcore-payments/ (accessed Nov.~12, 2026).

\bibitem{b26} E. Debenedetti et al., ``AgentDojo: A dynamic environment to evaluate prompt injection attacks and defenses for LLM agents,'' in Proc. NeurIPS Datasets and Benchmarks Track, 2024.

\bibitem{b27} T. South et al., ``Authenticated delegation and authorized AI agents,'' arXiv:2501.09674, Jan.~2025.

\bibitem{b28} A. Chan et al., ``Visibility into AI agents,'' in Proc. ACM Conf. Fairness, Accountability, and Transparency (FAccT), 2024, pp.~958--973, doi: 10.1145/3630106.3658948.

\bibitem{b29} S. Rose, O. Borchert, S. Mitchell, and S. Connelly, ``Zero Trust Architecture,'' NIST SP 800-207, Aug.~2020.

\bibitem{b30} D. R. MacIver, Z. Hatfield-Dodds et al., ``Hypothesis: A new approach to property-based testing,'' Journal of Open Source Software, vol.~4, no. 43, p.~1891, 2019, doi: 10.21105/joss.01891.

\bibitem{b31} S. Dashrath, ``kya-reference: Know Your Agent reference implementation,'' v1.0.2, GitHub, 2026. {[}Online{]}. Available: https://github.com/SwapnaliDashrath/kya-reference (accessed Nov.~12, 2026).
\end{thebibliography}
\end{document}
